\chapter{MÔ TẢ NGHIỆP VỤ}

\section{Tổng quan các nghiệp vụ}

Hệ thống quản lý kho hàng đảm nhận việc kiểm soát toàn bộ vòng đời của hàng hóa từ khi nhập vào kho cho đến khi xuất ra khỏi kho. Các nghiệp vụ chính bao gồm:

\begin{itemize}[leftmargin=1.5cm]
    \item \textbf{Nhập kho}: Tiếp nhận hàng hóa từ nhà cung cấp, kiểm tra chất lượng, số lượng và ghi nhận vào hệ thống.
    \item \textbf{Xuất kho}: Lập phiếu xuất kho cho khách hàng hoặc sản xuất, cập nhật số lượng tồn kho.
    \item \textbf{Kiểm kê}: Định kỳ kiểm tra thực tế hàng tồn kho so với sổ sách, điều chỉnh chênh lệch nếu có.
    \item \textbf{Quản lý hàng hóa}: Cập nhật thông tin hàng hóa, nhóm hàng, đơn vị tính, giá nhập/xuất.
    \item \textbf{Quản lý nhà cung cấp và khách hàng}: Lưu trữ, tra cứu và cập nhật thông tin đối tác.
    \item \textbf{Quản lý kho}: Tạo kho, cập nhật thông tin kho, phân bổ hàng hóa giữa các kho.
    \item \textbf{Báo cáo}: Tổng hợp số liệu nhập xuất tồn, hàng tồn lâu ngày, hàng sắp hết.
\end{itemize}

\section{Nghiệp vụ 1: Quản lý tài khoản người dùng (UC01)}

Quản trị viên tạo và quản lý tài khoản người dùng, phân quyền truy cập theo vai trò. Nghiệp vụ đảm bảo chỉ người dùng được ủy quyền mới thực hiện các thao tác trên hệ thống.

\section{Nghiệp vụ 2: Đăng nhập / Đăng xuất (UC02)}

Người dùng truy cập hệ thống bằng tên đăng nhập và mật khẩu. Hệ thống kiểm tra thông tin, xác thực và chuyển hướng đến giao diện phù hợp với quyền hạn.

\section{Nghiệp vụ 3: Quản lý nhà cung cấp (UC03)}

Nhân viên kho hoặc quản lý kho thêm, cập nhật và tra cứu thông tin nhà cung cấp (tên, địa chỉ, liên hệ). Thông tin này được sử dụng khi lập phiếu nhập kho.

\section{Nghiệp vụ 4: Quản lý khách hàng (UC04)}

Tương tự như quản lý nhà cung cấp, nghiệp vụ này lưu trữ thông tin khách hàng để phục vụ lập phiếu xuất kho và theo dõi giao dịch.

\section{Nghiệp vụ 5: Quản lý hàng hóa (UC05)}

Nhân viên kho nhập thông tin hàng hóa (mã, tên, nhóm hàng, đơn vị tính, giá), phân loại hàng hóa và cập nhật trạng thái sử dụng.

\section{Nghiệp vụ 6: Tìm kiếm hàng hóa (UC06)}

Người dùng tìm kiếm hàng hóa theo mã, tên, nhóm hàng hoặc kho. Hệ thống trả về thông tin hàng hóa và số lượng tồn kho hiện tại.

\section{Nghiệp vụ 7: Quản lý kho (UC07)}

Quản lý kho tạo và cập nhật thông tin kho (mã kho, tên kho, địa điểm, người phụ trách). Một hệ thống có thể có nhiều kho khác nhau.

\section{Nghiệp vụ 8: Tìm kiếm kho (UC08)}

Người dùng tra cứu thông tin kho, danh sách hàng hóa trong kho hoặc tồn kho theo từng kho.

\section{Nghiệp vụ 9: Lập phiếu nhập kho (UC09)}

Nhân viên kho lập phiếu nhập khi hàng hóa được chuyển đến từ nhà cung cấp. Phiếu nhập bao gồm thông tin nhà cung cấp, kho nhập, ngày nhập và danh sách hàng hóa nhập. Sau khi được duyệt, hệ thống tự động cập nhật tồn kho.

\section{Nghiệp vụ 10: Lập phiếu xuất kho (UC10)}

Nhân viên kho lập phiếu xuất khi có yêu cầu xuất hàng cho khách hàng hoặc sản xuất nội bộ. Hệ thống kiểm tra tồn kho trước khi cho phép xuất và cập nhật số lượng tồn sau khi duyệt.

\section{Nghiệp vụ 11: Quản lý đơn đặt hàng (UC11)}

Khách hàng hoặc bộ phận bán hàng tạo đơn đặt hàng. Hệ thống theo dõi trạng thái đơn hàng và tạo phiếu xuất kho khi đơn được duyệt.

\section{Nghiệp vụ 12: Quản lý thanh toán (UC12)}

Ghi nhận các khoản thanh toán liên quan đến phiếu nhập/xuất, cập nhật công nợ và tạo báo cáo tài chính tồn kho.

\section{Nghiệp vụ 13: Xem báo cáo thống kê (UC13)}

Quản lý kho và quản trị viên xem các báo cáo: nhập xuất tồn theo thời gian, hàng bán chạy, hàng tồn lâu ngày, giá trị tồn kho, cảnh báo tồn kho thấp/cao.

\section{Sơ đồ hoạt động (Activity Diagram)}

\subsection{Quy trình nhập kho}

\input{diagrams/activity_nhapkho}

\subsection{Quy trình xuất kho}

\begin{figure}[H]
\centering
\begin{tikzpicture}[
    startstop/.style={rectangle, rounded corners, draw, fill=red!15, minimum width=3cm, minimum height=0.8cm, font=\small, align=center},
    process/.style={rectangle, draw, fill=blue!10, minimum width=4cm, minimum height=0.8cm, font=\small, align=center},
    decision/.style={diamond, draw, fill=yellow!15, aspect=2, font=\small, align=center},
    arrow/.style={-{Stealth[scale=0.8]}, thick}
]

\node[startstop] (start) at (0,6) {Bắt đầu};
\node[process] (step1) at (0,4.8) {Nhân viên kho lập phiếu xuất kho};
\node[process] (step2) at (0,3.6) {Nhập thông tin khách hàng, kho, ngày xuất};
\node[process] (step3) at (0,2.4) {Thêm chi tiết hàng hóa xuất};
\node[decision] (dec1) at (0,1) {Tồn kho đủ?};
\node[process] (step4) at (0,-0.5) {Gửi phiếu xuất để duyệt};
\node[decision] (dec2) at (0,-2) {Quản lý duyệt?};
\node[process] (step5) at (0,-3.5) {Cập nhật tồn kho};
\node[process] (step6) at (0,-4.7) {In phiếu xuất};
\node[startstop] (end) at (0,-5.9) {Kết thúc};

\node[process] (err1) at (4,1) {Thông báo không đủ hàng};
\node[process] (err2) at (4,-2) {Thông báo từ chối};

\draw[arrow] (start) -- (step1);
\draw[arrow] (step1) -- (step2);
\draw[arrow] (step2) -- (step3);
\draw[arrow] (step3) -- (dec1);
\draw[arrow] (dec1) -- node[right, font=\scriptsize] {Có} (step4);
\draw[arrow] (dec1) -- node[above, font=\scriptsize] {Không} (err1);
\draw[arrow] (err1) |- (step2);
\draw[arrow] (step4) -- (dec2);
\draw[arrow] (dec2) -- node[right, font=\scriptsize] {Có} (step5);
\draw[arrow] (dec2) -- node[above, font=\scriptsize] {Không} (err2);
\draw[arrow] (err2) |- (step1);
\draw[arrow] (step5) -- (step6);
\draw[arrow] (step6) -- (end);

\end{tikzpicture}
\caption{Sơ đồ hoạt động quy trình xuất kho}
\label{fig:activity-xuatkho}
\end{figure}


\subsection{Quy trình kiểm kê tồn kho}

\begin{figure}[H]
\centering
\begin{tikzpicture}[
    startstop/.style={rectangle, rounded corners, draw, fill=red!15, minimum width=3cm, minimum height=0.8cm, font=\small, align=center},
    process/.style={rectangle, draw, fill=blue!10, minimum width=4cm, minimum height=0.8cm, font=\small, align=center},
    decision/.style={diamond, draw, fill=yellow!15, aspect=2, font=\small, align=center},
    arrow/.style={-{Stealth[scale=0.8]}, thick}
]

\node[startstop] (start) at (0,6) {Bắt đầu};
\node[process] (step1) at (0,4.8) {Nhân viên kho tạo phiếu kiểm kê};
\node[process] (step2) at (0,3.6) {Chọn kho và ngày kiểm kê};
\node[process] (step3) at (0,2.4) {Kiểm đếm thực tế từng hàng hóa};
\node[decision] (dec1) at (0,1) {Có chênh lệch?};
\node[process] (step4) at (0,-0.5) {Ghi nhận chênh lệch};
\node[process] (step5) at (0,-1.7) {Điều chỉnh tồn kho};
\node[process] (step6) at (0,-2.9) {Lập biên bản kiểm kê};
\node[startstop] (end) at (0,-4.1) {Kết thúc};

\node[process] (ok) at (4,1) {Khớp sổ sách};

\draw[arrow] (start) -- (step1);
\draw[arrow] (step1) -- (step2);
\draw[arrow] (step2) -- (step3);
\draw[arrow] (step3) -- (dec1);
\draw[arrow] (dec1) -- node[right, font=\scriptsize] {Có} (step4);
\draw[arrow] (dec1) -- node[above, font=\scriptsize] {Không} (ok);
\draw[arrow] (ok) |- (step6);
\draw[arrow] (step4) -- (step5);
\draw[arrow] (step5) -- (step6);
\draw[arrow] (step6) -- (end);

\end{tikzpicture}
\caption{Sơ đồ hoạt động quy trình kiểm kê tồn kho}
\label{fig:activity-kiemke}
\end{figure}


\section{Kết luận chương}

Chương 2 đã mô tả chi tiết 13 nghiệp vụ chính, đồng thời trình bày các sơ đồ hoạt động cho quy trình nhập kho, xuất kho và kiểm kê. Các nghiệp vụ và quy trình này là cơ sở để xác định tác nhân, use case và thiết kế các sơ đồ UML trong các chương tiếp theo.
